%-------------------------------------------------------------------------------
%	SECTION TITLE
%-------------------------------------------------------------------------------
\cvsection{Education}

%-------------------------------------------------------------------------------
%	CONTENT
%-------------------------------------------------------------------------------
\begin{cventries}
%---------------------------------------------------------
  \cventry
    {Instituto Superior Técnico} % Institution
    {M.Sc. Computer Science and Engineering} % Degree
    {Lisbon, Portugal} % Location
    {Sep 2023 - Nov 2025} % Date(s)
    {
      \begin{cvitems}
        \item {Specialized in \cvref{https://fenix.tecnico.ulisboa.pt/cursos/meic-t/curriculo?year=3380143556919297\#:~:text=168.0-,Inteligência Artificial\&text=168.0-,Artificial Intelligence}{\textbf{Artificial Intelligence}} with emphasis in \cvref{https://fenix.tecnico.ulisboa.pt/cursos/meic-t/curriculo?year=3380143556919297\#:~:text=Engenharia e Ciência de Dados\&text=Engineering and Data Science}{\textbf{Engineering and Data Science}}, combining theory with practical applications in modeling and data-driven problem solving.}
        \item {\textbf{Relevant Coursework}: \cvref{https://fenix.tecnico.ulisboa.pt/disciplinas/OA22/2024-2025/1-semestre}{Optimization Algorithms}, \cvref{https://fenix.tecnico.ulisboa.pt/disciplinas/AP-Dei/2023-2024/1-semestre}{Deep Learning}, \cvref{https://fenix.tecnico.ulisboa.pt/disciplinas/AVExe23/2023-2024/2-semestre}{Cloud Computing and Virtualization}, \cvref{https://fenix.tecnico.ulisboa.pt/disciplinas/AAut11/2024-2025/1-semestre}{Machine Learning}.}
        \item {Thesis: \cvref{https://fenix.tecnico.ulisboa.pt/cursos/meic-t/dissertacao/846778572214917}{\textbf{Maximum Variance Unfolding on Disjoint Manifolds (MVU-DM)}}.
          \begin{cvitems}
            \item {Developed a novel extension to the MVU algorithm that enables \textbf{nonlinear dimensionality reduction} on disconnected datasets.}
            \item {Engineered a decomposition-based approach that \textbf{allows for parallel processing}, improving scalability for large-scale datasets that were previously computationally infeasible.}
            \item {Preserved local geometric structure while achieving embeddings comparable to the standard method.}
          \end{cvitems}
        }
        \href{https://fenix.tecnico.ulisboa.pt/cursos/meic-t/curriculo?year=3380143556919297\#:~:text=Engenharia\%20e\%20Ci\%C3\%AAncia\%20de\%20Dados\&text=Data\%20Science}{Engenharia e Ciência de Dados / Data Science}
      \end{cvitems}
    }
    {
      Source notes, not rendered in the PDF.

      Context:
      - Master of Science in Computer Science and Engineering.
      - Final grade: 15/20.
      - Self-directed focus on Artificial Intelligence and Data Science.

      Main academic themes:
      - Deep learning and how LLMs work.
      - Statistics and calculus behind deep learning.
      - Planning, learning, and intelligent decision making.
      - Multi-agent systems.
      - Optimization and algorithmic foundations.

      Coursework groups:
      - Artificial Intelligence: Deep Learning, Autonomous Agents and Multi-Agent Systems, Planning, Learning and Intelligent Decision Making, Data Science.
      - Data Science: Data Analysis and Integration, Data and Information Systems Management, Cloud Computing and Virtualization.
      - Additional focus: Optimization Algorithms, Machine Learning.

      Thesis:
      - Maximum Variance Unfolding on Disjoint Manifolds (MVU-DM).
      - Decomposition-based extension to MVU for nonlinear dimensionality reduction on disconnected datasets.
      - Parallelizable design that improved scalability for large-scale datasets.
      - Comparable embeddings while preserving local geometric structure.

      Technologies used across coursework and thesis:
      - Python (NumPy, Pandas, Matplotlib, TensorFlow, Keras, Scikit-learn).
      - SQL (PostgreSQL, pgAdmin, SQL Server Management), MySQL, Microsoft SQL Server.
      - Pentaho (PDI, Server, Schema Workbench, Report Designer), DataCleaner, Saiku Analytics.
      - CmapTools, ArchiMate, Camunda Modeler, Enterprise Architecture (Sparx).

      Skills and signals:
      - Technical writing and clear presentation.
      - Working with large datasets and analysis tools.
      - Developing and evaluating machine learning models.

      End of source notes.
    }
%---------------------------------------------------------
  \cventry
    {Instituto Superior Técnico} % Institution
    {B.Sc. Computer Science and Engineering} % Degree
    {Lisbon, Portugal} % Location
    {Sep 2020 - Jul 2023} % Date(s)
    {
      \begin{cvitems}
        \item {\textbf{Built} a strong foundation in \cvref{https://fenix.tecnico.ulisboa.pt/disciplinas/ASA2/2021-2022/1-semestre}{Analysis and Synthesis of Algorithms}, \cvref{https://fenix.tecnico.ulisboa.pt/disciplinas/SDis/2022-2023/2-semestre}{Distributed Systems}, \cvref{https://fenix.tecnico.ulisboa.pt/disciplinas/SO2/2021-2022/1-semestre}{Operating Systems}, and \cvref{https://fenix.tecnico.ulisboa.pt/disciplinas/BD2/2021-2022/2-semestre}{Database Theory}.}
        \item {\textbf{Worked} on full-stack and systems-level projects with emphasis on \textbf{memory management} in \textbf{C/C++} and relational schema design in \textbf{SQL}.}
      \end{cvitems}
    }
    {
      Source notes, not rendered in the PDF.

      Context:
      - Bachelor of Science in Computer Science and Engineering.
      - Final grade: 13/20.
      - All courses were mandatory.
      - No scholarships or distinctions.
      - Participated in the annual snow trip to Andorra.

      Main academic themes:
      - Algorithms, distributed systems, operating systems, database theory, and systems thinking.
      - Practical experience with full-stack and systems-level projects.
      - Memory management in C/C++ and relational schema design in SQL.

      Coursework highlights:
      - Database Theory: relational model, schema design, implementation, and transaction management.
      - Analysis and Synthesis of Algorithms: algorithmic analysis, data structures, complexity, and design.
      - Distributed Systems: scalability, fault tolerance, security, remote calls, and distributed naming.
      - Operating Systems: system-level programming and OS internals.

      Technologies used:
      - C, C++, Java, Python (NumPy, Pandas).
      - SQL (PostgreSQL), pgAdmin, Git, Linux, Bash Scripting, RESTful APIs, LaTeX.

      Skills and signals:
      - Collaboration on complex technical projects.
      - Persistence in problem-solving.

      End of source notes.
    }
%---------------------------------------------------------
\end{cventries}
